% =============================================
% BYLAWS CONTENT FILE
% Edit ONLY this file for content changes.
% Use normal LaTeX enumerate environments for all numbered/lettered lists.
% The formatting is handled automatically in main.tex
% =============================================

\begin{enumerate}
    \item These shall be the procedures that govern the hearings to be held by all Institute Hearing Boards and the Review Board. *
    \item All appeals to an Institute Hearing Board, or the Review Board shall be made through the Office of the Dean of Students. All appeal requests shall be made in writing and can only be based on one or more of the three Institute grounds for appeal and must specifically state the reason and rationale for appeal.
    \item Prior to a hearing both parties shall independently provide, to the Board Chair, all facts and evidence they wish to present at the hearing, no later than the deadline determined by the Chair. This evidence will be in writing and will contain such items as incident reports, Public Safety reports, judicial inquiry notes, etc. The evidence will also include any and all witnesses the parties wish to have testify. Following said witness requests, those witnesses determined to be relevant and vital to the hearing, will be approved. In order to be considered for approval, this request must include the specific relevancy of a witness, sufficient justification to warrant them testifying at the hearing, and the full name, title, current address and phone number for each potential witness. If any of the previously mentioned information is not provided, is past the deadline, or if the witness is determined to be irrelevant, the witness and his/her testimony will not be admissible in the hearing. Upon receipt of the evidence from both parties, the Board Chair will review the information for relevancy and compile the approved evidence from both parties into a single evidence packet. A copy of this packet will then be distributed to both parties, all board members, and the Senior Judicial Administrator or designee. This evidence exchange, between both parties, will ideally take place at the earliest availability of the evidence, but not less than 48 hours prior to a hearing. This report many include a statement of the student's/group's past disciplinary record(s) obtained from the Office of the Dean of Students. This portion of the report is used for sanctioning purposes, as appropriate.
    \item A lower board may refer a case directly to the Review Board. In doing so, the lower board shall provide a written statement delineating its reasons. The Review Board may similarly refer a case to the President.
    \item The Chair or appointee shall set the date, time and place for the hearing and shall notify all other board members and parties. This hearing will be held at least three (3) Institute business days but not more than ten (10) Institute business days after receipt of the hearing request (Rensselaer reserves the right to suspend the judicial process near or during holidays, breaks, summer, etc.). The chair shall give written notice of the hearing to the parties involved. Board hearings will be held weekdays, between the hours of 9:00m and 7:00pm and for a duration not to exceed 2.5 hours per day/segment. Exceptions to this policy will be made only in extreme circumstances and at the sole discretion of the Board Chair or designee and the Senior Judicial Administrator.
    \item If either party finds it necessary to postpone a board hearing, they must contact the Chair of the Board, the Vice-Chair, or the Senior Judicial Administrator, not less than 24 hours before the scheduled time of the hearing. At the discretion of the Chair, the case may be rescheduled. If the 24- hour condition is not met, the Board may hear the case at the originally scheduled time or uphold the original decision in the case. Should the Board find itself unable to convene for the scheduled case, it may notify the parties involved and reschedule the case for the earliest possible time both parties and the Board can be present. Postponements will only be considered in extreme and unforeseeable cases (i.e., a death in the family of one of the participants).
    \item All parties shall be furnished with a copy of these procedures and the appropriate bylaws prior to the hearing, at a meeting with the Chair. At this meeting, the party may petition the Chair (in writing) to disqualify a student board member form a particular case with good cause. If said member does not then disqualify himself or herself, the Board may, by a $2/3$ vote of its membership, disqualify that member from that particular case for the reasons stated in the petition if the reasons are determined to be valid.
    \item The hearing will be closed except for persons whose presence is authorized by the Board Chair. However, victims (defined as individual victims of physical violence-including sexual offenses) shall be allowed to be present throughout the duration of the hearing, if that individual requests and so chooses to do so.
    \item The Chair shall rule on questions of relevancy and procedure prior to the hearing. The Chair will also rule on questions of relevancy and procedure during the hearing, but can be overruled by a $2/3$ votes of the Board, with good cause.
    \item The Senior Judicial Administrator shall be present throughout the duration of a board's hearing/evidence review, deliberation, and decision making to provide clarification, advise on questions of relevancy, issues of procedure and structure, and to respond to questions of process or other inquiries.
    \item At the hearing, both parties shall have the right to be assisted by a Student Judicial Advisor (if requested) and the right to testify, to produce evidence (the evidence packet), and to cross-examine witnesses. This advisor must be a current member of the Rensselaer student community and officially designated as a Student Judicial Advisor. The advisor shall not take part in the proceedings except to advise the accused, the witness, the complainant, etc.
    \item At the hearing, the accused shall be introduced to all board members, and shall be informed of the appeal procedures.
    \item The hearing shall proceed as follows:
    \begin{enumerate}
        \item The party bearing the burden of proof shall make his/her opening statement.
        \item The opposing party shall follow with his/her opening statement.
        \item Any witnesses subpoenaed by the party bearing the burden of proof shall be called first, then any witnesses called by the opposing party, and finally any witnesses called by the Board. All witnesses shall be questioned first by the person calling them, then by the opposing person, and finally by the members of the Board.
        \item Following the testimony of witnesses, the Board may question either party, and the two parties may question each other.
        \item Following the period of questioning, the party bearing the burden of proof shall make his/her closing statement.
        \item The opposing party will then make his/her closing statement.
        \item After both parties have presented their closing statements, the hearing will be closed for deliberations.
    \end{enumerate}
    \item The Chair shall provide the Senior Judicial Administrator with the decision in writing within two Institute business days of the decision.
    \item Within three (3) Institute business days after the close of the Board hearing, the Senior Judicial Administrator shall inform (at a meeting or in writing only) the accused of the Board's decision and recommendation(s).
    \item The Board shall make a tape recording of the hearing. The decision of the Board shall be kept on file within the restrictions of confidentiality required by Institute policy.
    \item The Chair, in consultation with the Senior Judicial Administrator, and with the consent of both parties, can suspend the above procedures for a particular case with good cause.
    \item A Board may issue a temporary injunction to stop an action from continuing or to prevent a potential action from occurring. However, this injunction may be changed as a result of a hearing by the respective Board.
    \item If an approved witness (see item \#3) is a current Rensselaer student, the board can subpoena him/her to testify. A subpoena may be issued, by the Chair, upon appropriate request from either party, to the individuals who are members of the Rensselaer student community. All parties are subject to the witness submission deadline given to them by the board Chair (witness submissions that are incomplete or after the determined deadline are not admissible). A subpoena cannot be issued any less than 72 hours prior to the given case hearing. If a relevant and approved witness is not a current Rensselaer student, it is the responsibility of each party to request that he/she serve as their witness and to notify him/her of the hearing date and time. As a general rule, the business of the board will not be delayed due to the schedule of a witness or his/her unavailability. Any student witness testifying may be assisted by an advisor. This advisor must be an officially designated Student Judicial Advisor.
    \item The Judicial Board may propose amendments to these procedures. Before any amendment becomes effective, it shall be approved by the President of the Institute not less than ten days nor more than 30 days after it is proposed. The Judicial Board must inform the Student Senate and Faculty Senate when changes have been made to the hearing procedures.
    \item The standard of proof in judicial proceedings is a preponderance of the evidence (51\%).
\end{enumerate}

* The Review Board conducts a review of the written evidence (including the evidence packet compiled for the lower board hearing and the transcript of said hearing) to determine it's verdict and/or sanctions(s). Therefore, with the exception of needing to decide a case within three to ten business days of receiving the appeal request, the Review Board may only be subject to certain hearing procedures in the unlikely event that it holds an actual hearing to decide on an appeal.